\documentclass[11pt]{article}
\usepackage[a4paper,margin=27mm]{geometry}
\usepackage{amsmath,amssymb,amsthm,microtype}
\usepackage[T1]{fontenc}
\usepackage[colorlinks=true,linkcolor=blue,citecolor=blue,urlcolor=blue]{hyperref}
\newtheorem{theorem}{Theorem}
\newtheorem{lemma}{Lemma}
\newcommand{\E}{\mathbb E}
\newcommand{\Prob}{\mathbb P}
\newcommand{\Z}{\mathbb Z}
\newcommand{\one}{\mathbf 1}
\newcommand{\negp}[1]{(#1)_{-}}
\hypersetup{
 pdftitle={Near-one oriented-percolation flow from Bernoulli first-passage estimates},
 pdfauthor={Alec Kriebel},
 pdfsubject={An attributed corollary of BGT's 2021 estimates for the Aldous-Lenderman flow asymptotic},
 pdfkeywords={oriented percolation, maximum flow, Bernoulli first-passage percolation, strict chains}
}
\title{Near-one oriented-percolation flow\\from Bernoulli first-passage estimates}
\author{Alec Kriebel\\\small Independent researcher\\
\small\href{https://orcid.org/0009-0001-9320-500X}{ORCID: 0009-0001-9320-500X}}
\date{16 October 2026 \quad (version 1.0.0)}
\begin{document}
\maketitle

\begin{abstract}
For the square-lattice oriented Bernoulli edge-flow density, we derive
$1-v(p)\sim\sqrt{2(1-p)}$ as $p\uparrow1$, the asymptotic posed by
Aldous and Lenderman. The argument is an explicit corollary of
Basdevant--Gou\'er\'e--Th\'eret's 2021 strict-chain scaling and
weighted-list estimate. The directed cut dual has zero-cost backward
steps, so their ordinary first-passage theorem cannot be substituted
directly. Grouping the two outgoing rare-edge marks into independent
site marks gives the required event inclusion into their already
evaluated counting sum. We supply this comparison, the finite
directed-cut reduction and the order of limits. The asymptotic is
presented as an attributed application of older estimates; no claim of
novelty, historical firstness or present-day openness is made.
\end{abstract}

\section{The question and the precise model}

Aldous's problem page \cite{Aldous} describes the conjecture, formulated
with Jason Lenderman in 2004, and the difficulty of justifying its
Poisson--Hammersley heuristic. We consider exactly that edge-flow model.
The vertices are $\{0,\ldots,n-1\}^2$, $n\ge2$, with edges pointing north
or east. Independently, each edge has capacity one with probability $p$
and capacity zero otherwise. The source vertices are the left and
bottom boundary vertices except $(0,n-1)$ and $(n-1,0)$; the sinks are
the top and right boundaries. Let $V_n$ be the maximum flow between
these sets. Integrality identifies it with the maximum number of
edge-disjoint open directed paths; shared vertices are permitted.

The source states the deterministic $L^1$ limit
$V_n/(2n)\longrightarrow v(p)$. We use this stated existence result
with its normalization. In fact the expectation limit used below also
follows directly from diagonal superadditivity after the finite
reduction. No heuristic identity from the problem page is an input.

\begin{theorem}\label{main}
In this model,
\[
 \lim_{q\downarrow0}\frac{1-v(1-q)}{\sqrt{2q}}=1.
\]
Here the finite-volume limit defining $v$ is taken at each fixed $q>0$
before $q$ tends to zero.
\end{theorem}

The substantial older input is the numerical weighted-list bound
inside the proof of Basdevant--Gou\'er\'e--Th\'eret (BGT)
\cite[\S3.3, equation (33)]{BGT}, together with their strict-site score
scaling. All equation numbers for that input refer to their public
2021 arXiv v1; publication metadata identifies a 2022 journal version.
We do not assert a line-by-line identity of the two editions.

\section{The finite directed-cut reduction}

Put $N=n-2$ and
\[
 e_L=((0,n-2),(0,n-1)),\qquad
 e_B=((n-2,0),(n-1,0)).
\]
Write $c_L,c_B$ for their capacities. Reflect the vertical coordinate
of the bounded-face dual, obtaining $Q_N=\{0,\ldots,N\}^2$. A dual
east or north traversal costs the capacity of the primal edge it
crosses; a reverse traversal costs zero. For example, a dual east
edge starting at $(x,y)$ crosses
$((x+1,N-y),(x+1,N+1-y))$, and a dual north edge crosses
$((x,N-y),(x+1,N-y))$. The correspondence is injective on unoriented
dual edges. Let $T_N$ be the minimum dual cost from $(0,0)$ to
$(N,N)$, with $T_0=0$.

\begin{lemma}\label{dual}
For every nonnegative primal capacity assignment,
\begin{equation}\label{identity}
 V_n=c_L+c_B+T_N.
\end{equation}
\end{lemma}

\begin{proof}
Consider any cut whose source side contains all prescribed sources
and no sink. Draw the crossing of every primal edge with differently
labelled endpoints, orienting the interface with the source side to
its right in the physical embedding. Include exterior endpoints at
$e_L,e_B$. At each bounded face the numbers of incoming and outgoing
crossings agree: the two types of binary-label transitions alternate
around its four corners. This includes a checkerboard face with two
incoming and two outgoing crossings. The endpoint at $e_L$ has one
outgoing crossing, that at $e_B$ one incoming crossing, and no other
boundary edge changes labels.

There is a directed interface route between these endpoints.
Otherwise the set reachable from the first endpoint has no outgoing
arc, contradicting its total outdegree-minus-indegree balance.
Extract a path and erase loops. Forward dual crossings represent
outward primal cut edges and reverse crossings inward edges; only the
former contribute to the directed cut. Nonnegativity therefore gives
cut cost at least $c_L+c_B+T_N$.

Conversely, a vertex-simple dual corner path, extended through
$e_L,e_B$, is a simple crosscut of the primal square meeting no
primal vertex. Planar separation gives a feasible primal cut on the
side containing the source boundary arc. Its outward crossings are
exactly the forward dual crossings and the two forced edges.
Thus its cost is the cost in \eqref{identity}. A minimizing dual walk
can be made simple by deleting nonnegative-cost cycles. Apply
max-flow/min-cut. For $n=2$ there is no internal dual step and the same
argument gives $V_2=c_L+c_B$.
\end{proof}

Set $q=1-p$ and call a forward dual edge \emph{marked} when its cost
is zero. These marks are independent Bernoulli$(q)$. For a simple
corner path let $k$ count marked forward steps and $b$ backward
steps. Its forward count is $2N+b$, so its cost is $2N+b-k$.
Consequently
\begin{equation}\label{defect}
 D_N:=2N-T_N=\max_P(k(P)-b(P)),\qquad 0\le D_N\le2N.
\end{equation}
The maximum can be restricted to simple paths. Equivalently every
cycle has nonpositive reward, since it has equally many forward and
backward steps.
Writing $g(q)=1-v(1-q)$, bounded boundary terms in \eqref{identity}
give
\begin{equation}\label{density}
 g(q)=\lim_{N\to\infty}\frac{\E D_N}{2N}.
\end{equation}
The limit of these expectations exists even without the source's
$L^1$ premise: concatenating optimal paths in consecutive diagonal
squares makes $\E D_N$ superadditive, bounded above by $2N$.
The elementary superadditive limit lemma applies.

\section{The older strict-chain and list estimates}

Extend the forward marks independently outside $Q_N$. Declare a site
$x$ open when at least one of its two outgoing forward edges is marked.
Different sites use disjoint pairs of edge coordinates. Thus
these site indicators are independent with parameter
\begin{equation}\label{r}
 r=1-(1-q)^2=2q-q^2.
\end{equation}
Let $A_r(M)$ be the maximum length of a coordinatewise strictly
increasing open-site chain in $[0,M)^2\cap\Z^2$, allowing an open
origin. This is BGT's site score $S_r^D((M,M))$ in dimensions
$d_1=2,d_2=0$. Their Lemmas 5, 10 and 13, with the planar
Poisson-chain constant identified in their \S3.4, give
\begin{equation}\label{chain}
 \sigma_r=\lim_{M\to\infty}\frac{\E A_r(M)}{M},
 \qquad \lim_{r\downarrow0}\frac{\sigma_r}{\sqrt r}=2.
\end{equation}

We state precisely the other consequence of their counting
calculation that will be used. This concerns an integer list
$w_0=0,w_1,\ldots,w_p,w_{p+1}=(M,M)$, with the real intermediate
sites distinct, not equal to zero, and strictly below $(M,M)$ in
each coordinate. Starting with index zero, repeatedly select the
first subsequent site strictly larger in both coordinates. The
artificial terminal guarantees termination. Write $a$ for the number
of selected gaps and $\ell_j$ for the number of unselected sites in
gap $j$, $0\le j<a$. Thus
\[
 L=\sum_j\ell_j,\qquad p=a-1+L.
\]
Let $c(j,0),\ldots,c(j,\ell_j)$ be the selected starting anchor and
the intervening sites, and let $c(j+1,0)$ denote the next anchor.
Set
\[
 h(j)=c(j+1,0)-c(j,0),\qquad
 u(j,i)=c(j,i)-c(j,i-1).
\]
If $\ell_j>0$, choose a coordinate $d_j$ in which
$c(j,\ell_j)_{d_j}\le c(j,0)_{d_j}$; such a coordinate exists by
the first-subsequent selection rule. For an empty gap choose either
coordinate. Let $d'_j$ be the other coordinate and define
\begin{equation}\label{X}
 X=\sum_j\sum_{i=1}^{\ell_j}|u(j,i)_{d_j}|
   +\sum_j\sum_{i=1}^{\ell_j}\negp{u(j,i)_{d'_j}}.
\end{equation}
Here $t_-=\max(-t,0)$. The second term is the \emph{sum of individual
negative parts}, as in BGT's (19) and Lemma 19; it is not the
negative part of the net displacement.

\begin{lemma}[BGT weighted-list consequence]\label{list}
Fix $\tau>0$. For all sufficiently small $r>0$, the sum of the
weights $r^p$ over such lists satisfying
\begin{align}
 c(j,\ell_j)_{d'_j}-c(j+1,0)_{d'_j}
       &\le2(\ell_j+1) \quad(0\le j<a),\label{item4}\\
 0&\le-X+1+2L-\tau M\sqrt r\label{item5}
\end{align}
is at most
\[
 \frac{C_\tau}{r}\exp(-M\sqrt r),
\]
for a finite constant $C_\tau$ independent of $M\ge2$.
\end{lemma}

\begin{proof}[Identification of the published estimate]
Use the $m=0$ numerical terms of BGT's (33), with
$d_1=2,d_2=0$. Their Items 1 and 3 hold by the telescoping positive
anchor increments and the choice of $d_j$; Item 2 has no coordinates.
Their remaining two predicates are \eqref{item4}--\eqref{item5}.
The weight is exactly $r^{a-1+L}=r^p$. This is a numerical sum over
encodings; it does not require a geodesic in the model originally
used to derive it.

For clarity, the specialization of their displayed $W$ bound has
the following factors. Put $K(t)=(1+e^{-t})/(1-e^{-t})$, choose
$\beta=2\gamma$, and set
\begin{align*}
 B_0(r)&=r e^2 e^{2\gamma\sqrt r}
                   K(\beta\sqrt r)K(\gamma\sqrt r),\\
 B_1(r)&=r e^{2\gamma\sqrt r+2\alpha} e^2 K(\alpha)
                   [K(\gamma\sqrt r)+K(\alpha-\gamma\sqrt r)].
\end{align*}
For fixed $a$, $(\ell_j)$ and $(d_j)$, their coordinate summation is
bounded by
\[
 r^{-1}e^{2+\alpha}
 e^{(2\beta-\alpha\tau)M\sqrt r}
 B_0(r)^a B_1(r)^L.
\]
We have written the gap as $\tau M\sqrt r$ in the numerical
predicate, thereby fixing its site-parameter normalization. No
identification of a bond-FPP event with a site event is used.

Since $K(t)\sim2/t$ as $t\downarrow0$,
$B_0(r)\to4e^2/(\beta\gamma)$ and $B_1(r)\to0$.
First choose $\gamma$ large, then $\alpha$ with
$\alpha\tau>2\beta+1$, then $r$ small so that
$\alpha-\gamma\sqrt r>0$ and
$2B_0(r)/(1-B_1(r))\le1/2$.
Summing the $2^a$ coordinate choices and the nonnegative
$\ell$-vectors gives the geometric series
\[
 \sum_{a\ge1}\left(\frac{2B_0(r)}{1-B_1(r)}\right)^a\le1.
 \]
 This gives the bound, for example with $C_\tau=e^{2+\alpha}$.
These are the already evaluated positive numerical terms of BGT's
calculation; enlarging the set of encodings only increases the
majorant.
\end{proof}

The meaning of $m=0$ is important. It prescribes no doubly-open
sites in that numerical weight. It does \emph{not} assert that the
actual path visits no site with two marked outgoing edges. For any
fixed list of distinct sites, requesting that its $p$ sites be open
has probability $r^p$, irrespective of which have two marked edges.

\section{The comparison for the asymmetric dual}

First, every strict open-site chain counted by $A_r(N)$ can be
visited by a monotone path in $Q_N$ that takes a marked outgoing
edge at each chain site. Both coordinates increase by at least one
between successive sites, so either outgoing edge fits before the
next site. Each chosen edge stays inside $Q_N$, including the final
connection. Thus, pathwise,
\begin{equation}\label{lower}
 D_N\ge A_r(N).
\end{equation}

\begin{lemma}\label{upper}
For every $\delta>0$, all sufficiently small fixed $q>0$ satisfy
\begin{equation}\label{tail}
 \Prob\{D_N>A_r(N+1)+\delta(N+1)\sqrt r\}
 \le \frac{C_{2\delta}}{r}\exp(-(N+1)\sqrt r),
 \qquad N\ge1.
\end{equation}
\end{lemma}

\begin{proof}
Choose a simple minimizing dual path $P$ in $Q_N$ and set $M=N+1$.
Every vertex of $P$ is strictly below the auxiliary terminal
$(M,M)$. List all open vertices visited by $P$, excluding its actual
terminal $(N,N)$ and treating the origin separately. Append the
auxiliary terminal and apply the extraction of Section 3. Let
$I_0$ indicate that the origin is open.

Every marked forward traversal has a distinct visited starting
vertex, because $P$ is simple. At most one such traversal leaves
each visited vertex, even if both outgoing edges are marked.
Therefore
\[
 k\le p+I_0,\qquad A_r(M)\ge a-1+I_0.
\]
Together with $D_N=k-b$ and $p=a-1+L$, this yields
\begin{equation}\label{cancellation}
 D_N-A_r(M)\le L-b.
\end{equation}
The origin indicator cancels exactly.

Consider a nonlast gap between two real selected anchors $s\prec t$.
Its actual subpath is a minimum-cost path between those anchors:
a cheaper replacement would give a cheaper global walk, and
nonnegative costs permit subsequent loop removal. A monotone
competitor can take a marked first edge if $s$ is open, since
both coordinate displacements to $t$ are at least one. Its cost is
at most $\|t-s\|_1-I_s$. The actual segment's cost is
$\|t-s\|_1+b_j-k_j$, with $k_j\le\ell_j+I_s$. Hence
\[
 b_j\le\ell_j.
\]
Any overshoot of the last intervening site beyond $t$ in the other
coordinate requires that many backward steps to return. This proves
\eqref{item4}, even with $\ell_j$ in place of $2(\ell_j+1)$.
For the last gap, all real sites are strictly below the auxiliary
terminal, so \eqref{item4} follows directly. No comparison path
outside $Q_N$ is needed. Artificially open boundary tails caused by
the independent extension only enlarge the upper chain score.

For a gap's distinguished coordinate, the prefix through its last
intervening site has nonpositive net displacement. Thus the sum of
absolute marked-site increments in this coordinate is at most twice
its actual backward variation. In the other coordinate the sum of
individual negative increments is at most the corresponding actual
backward variation. These prefixes are disjoint in path time across
gaps. Consequently the exact quantity \eqref{X} satisfies
\begin{equation}\label{variation}
 X\le2b.
\end{equation}
This factor two is retained, rather than identifying backward costs
in two different models.

On the event in \eqref{tail}, \eqref{cancellation}--\eqref{variation}
give
\[
 -X+2L>2\delta M\sqrt r.
\]
In particular \eqref{item5} holds with $\tau=2\delta$, including
its harmless additional $1$. The path therefore supplies a list of
distinct open sites meeting all predicates of Lemma \ref{list}.
For each fixed list its open-site probability is $r^p$; we may
ignore the origin's state. The union bound and that lemma prove
\eqref{tail}. There is no conditioning on a chosen geodesic and no
independence assumption between different lists.
\end{proof}

To finish Theorem \ref{main}, use $0\le D_N\le2N$ and $A_r\ge0$:
\[
 \E D_N\le\E A_r(N+1)+\delta(N+1)\sqrt r+
                 2N\,\Prob\{\text{event in \eqref{tail}}\}.
\]
First let $N\to\infty$ at a fixed sufficiently small $q$. Equations
\eqref{density}, \eqref{chain} and \eqref{lower} yield
\[
 \frac{\sigma_r}{2}\le g(q)
     \le\frac{\sigma_r}{2}+\frac{\delta\sqrt r}{2}.
\]
Next let $q\downarrow0$, using
$\sigma_r/\sqrt r\to2$ and $r/(2q)\to1$. Finally let $\delta\downarrow0$.
This proves the stated asymptotic with the required order of limits.

\section{Priority, scope and verification}

The entire asymptotic above is a corollary of BGT's older
strict-chain and weighted-list estimates after the explicit
directed-model event comparison. It is not a new central
rare-defect estimate. Our bounded priority checks did not locate
an earlier explicit statement of this particular flow corollary;
that fact neither proves first recognition nor establishes that the
historical question remained open. This note records the deduction
and its attribution without making either assertion.

The ordinary undirected FPP theorem alone is insufficient. In a
common box a path's ordinary deficit is $k_f+k_b-2b$, whereas the
present deficit is $k_f-b$. The present deficit minus the ordinary
deficit is $b-k_b\ge0$, unhelpful for the needed upper bound. The larger numerical
list estimate, rather than equality of the two models, closes the
comparison.

The supporting archive also preserves the original, longer
directed-cut counting and BK--Poisson proof and its audits. The
Poisson-chain input and upper tail are classical \cite{DZ}; so is
product-space disjoint occurrence \cite{BK}. Overlapping-region
decomposition by disjoint witnesses and endpoint skeletons also
predates this work, including Verges \cite[Lemma 2.5 and \S5.2]{Verges}.
No originality claim for those general mechanisms is made.
The present statement concerns the prescribed edge-capacity model;
it makes no assertion for vertex-disjoint paths, alternative
boundary conventions or other graph families.

Finite computations test directed duality, cut geometry, sequence
constraints and replay of the archived diagnostics. They supplement
the analytic proof and the cited estimates, and do not establish
the infinite-volume conclusion or novelty. The archive gives exact
source fingerprints and bounded-search limitations; third-party
publications and raw source corpora are not redistributed.

\paragraph{AI use and review status.}
AI tools were used extensively in mathematical development, source
review, writing, computation and adversarial checking. The recorded
adversarial reviews were performed by AI agents. This is an
unrefereed preprint; no completed independent human peer review or
formal proof-assistant verification is claimed. The author is
responsible for the claims.

\begin{thebibliography}{9}
\bibitem{Aldous}
D.~Aldous, \emph{A discrete Hammersley process as an extreme case of
oriented percolation flows}, problem page, conjecture formulated
with J.~Lenderman in 2004.
\url{https://www.stat.berkeley.edu/~aldous/Research/OP/hammersley_flow.html}.

\bibitem{BGT}
A.-L.~Basdevant, J.-B.~Gou\'er\'e and M.~Th\'eret,
First-order behavior of the time constant in Bernoulli first-passage
percolation, \emph{Ann. Appl. Probab.} \textbf{32} (2022), 4535--4567.
\href{https://doi.org/10.1214/22-AAP1795}{doi:10.1214/22-AAP1795}.
Operative primary version:
\href{https://arxiv.org/abs/2106.12266v1}{arXiv:2106.12266v1},
23 June 2021, especially \S\S2.2--2.5 and 3.3.

\bibitem{DZ}
J.-D.~Deuschel and O.~Zeitouni,
On increasing subsequences of i.i.d.\ samples,
\emph{Combin. Probab. Comput.} \textbf{8} (1999), 247--263.
\url{https://www.wisdom.weizmann.ac.il/~zeitouni/pdf/ccp.pdf}.

\bibitem{BK}
J.~van den Berg and H.~Kesten, Inequalities with applications to
percolation and reliability, \emph{J. Appl. Probab.}
\textbf{22} (1985), 556--569.
\href{https://doi.org/10.2307/3213860}{doi:10.2307/3213860}.

\bibitem{Verges}
J.~Verges, \emph{Concentration inequalities and large deviations for
continuous greedy animals and paths},
\href{https://arxiv.org/abs/2503.01484v1}{arXiv:2503.01484v1},
3 March 2025.
\end{thebibliography}
\end{document}
